\documentclass[11pt]{article}
\usepackage[utf8]{inputenc}
\usepackage[margin=1in]{geometry}
\usepackage{enumitem}

\begin{document}
\begin{titlepage}
    \centering
    \vspace*{2cm}
    {\Huge \textbf{Snap Shot objectives}} \\
    \vspace{0.5cm}
    {\Large \textbf{JPL Lunar Detection Pipeline}} \\
    \vfill
    \textbf{Prepared by:} \\
     Ricky Chao, George Malanche, Kevin Mendoza, Connor Moy \\
    \vspace{1cm}
    \textbf{Sponsored by:} NASA JPL \\
    \textbf{Date:} May 15, 2026
\end{titlepage}

\section*{Start Objective}
The starting objective of the JPL Lunar Detection Pipeline is to establish the core software infrastructure and data ingestion capabilities necessary for high-resolution lunar surface feature analysis. This includes the following key components:

\begin{itemize}
    \item \textbf{Data Ingestion Setup:} Establishing Phase 0 mechanisms for downloading LROC NAC products (.IMG and PDS4 .xml pairs) from the NASA PDS archive.
    \item \textbf{Preprocessing Framework:} Implementing Phase 1 imaging to convert raw products into analysis-ready GeoTIFF/PNG tiles.
    \item \textbf{SPICE Integration:} Configuring the NASA SPICE toolkit for initial camera geometry.
\end{itemize}

\section*{1st Checkpoint Objective}
The primary objective for Snapshot 2 is to implement the end-to-end annotation conversion and initial model training pipeline. This involves developing functional scripts for Phase 2 (annotation) and Phase 3 (train-yolo), including:
\begin{itemize}
    \item Converting PASCAL VOC XML annotations into standardized COCO and YOLO formats.
    \item Implementing dataset splitting (train/val/test) with reproducible shuffling.
    \item Executing baseline training runs on YOLOv11 to verify model convergence and basic feature identification.
\end{itemize}

\section*{2nd Checkpoint Objective}
The primary objective for Snapshot 3 is to integrate advanced detection engines and high-precision geolocation services. This phase focuses on the implementation of:
\begin{itemize}
    \item \textbf{Detectron2 Pipeline (Phase 4):} Training Mask R-CNN architectures to provide a high-accuracy detection alternative to the YOLO model.
    \item \textbf{Precision Geolocation (Phase 5):} Refining ray-tracing computations to map detections to precise lunar latitude and longitude coordinates.
\end{itemize}

\section*{Due Date Checkpoint}
The primary objective for the Due Date Checkpoint is to finalize the integrated pipeline and ensure full requirements validation. Final efforts will focus on:
\begin{itemize}
    \item \textbf{Error Handling and Logging:} Finalizing the Phase 5 inference engine with robust error reporting and rollback mechanisms for failed processing batches.
    \item \textbf{Analysis and Visualization:} Completing the visualization exporter to overlay geolocated bounding boxes on original lunar imagery.
    \item \textbf{Documentation and Verification:} Conducting end-to-end system testing to ensure the pipeline meets performance targets, such as achieving at least 90\% accuracy for crater detection.
\end{itemize}

\end{document}
